\section{Discussion}\label{sec:discussion}
Aggregate stressed ranking improves in both observed blocks, with a mean difference of approximately 0.1602 ROC-AUC units. A cleaner separation of many stressed sources can raise the average even if particular attack groups shift below the chosen threshold. The measured changes in attacks A01, A03 and A04 and the worsening on A05 and A06 show why aggregate ranking is insufficient for selecting an operating point. The latter two attacks are represented in training. Their deterioration reflects a within-benchmark robustness gap under the tested intervention.

The error effect depends on initialization. The A stressed FAKE miss reduction is larger for 9107, and the 9207 cluster interval includes zero. Clean B false alarms move in opposite directions, with the 9107 increase concentrated in two speakers. Changing the score distribution relative to 0.5 could contribute to these changes, even if discrimination improves. The results describe decisions under the registered threshold and do not evaluate threshold recalibration.

For many REAL sources, Noise-to-Opus is the only erroneous order. For FAKE sources, the reverse chain is often the only erroneous order. An upward shift in FAKE scores under one order can produce both patterns. The retrospective matched-component diagnostic for 9107 finds that nuisance adjustment does not establish an independent forensic target. Nonlinear codec response, preprocessing, and detector sensitivity may all contribute. Reproducing the direction of decision asymmetry in 9207 does not independently reproduce the nuisance analysis.

In practical evaluation, reporting both threshold-specific risks and ranking separates several failure modes that may have different costs. Speaker concentration also affects how a review workload might change, although no field workload or utility was measured. Conventional Opus order reversal, class-wise counts and a matched training intervention provide a focused complement to broader robustness benchmarks and codec-adaptation methods.

\subsection{Aggregate improvement and its distribution}
The matched comparison makes training-view policy the controlled factor, but does not isolate a single cause inside the fitted network. The channel model can learn features that distinguish more stress-processed pairs, move score distributions relative to the fixed decision boundary, or do both. Clean BCE selection may choose a different epoch in each arm because the fitted score distributions are different. The observed contrast therefore includes the interaction of augmentation with the same selection rule. It cannot determine which intermediate representation would be responsible, or how an alternate selection rule would change outcomes.

The adverse results for A05 and A06 matter in two distinct ways. First, these groups already contribute repeated stress errors before augmentation. Their denominators are small, so an individual source can affect several condition counts. Second, the registered model comparisons show additional errors distributed across several sources, especially in the 9207 block. The per-condition breakdown distinguishes an attack already difficult under Noise5 from additional failures concentrated in the ordered compositions. This supports a localization claim about the tested inputs. It does not show which characteristics of their synthesis pipelines caused the effect.

REAL false alarms can shift if a score displacement moves bona fide examples near 0.5 across the threshold. In B the two blocks move in opposite directions, and part of the 9107 increase concentrates in a small number of speakers. Bona fide cross-testing work makes the broader case for assessing bona fide diversity \cite{kwok}. Our speaker analysis addresses a narrower issue inside one fixed cohort. A different bona fide corpus, language, noise mixture or deployment prior could change both the baseline rates and the direction of the intervention effect.

\subsection{Signal-chain and engineering implications}
Published communication benchmarks such as ADD-C vary codec and packet-loss conditions \cite{addc}. Their breadth is useful for establishing stress coverage. The same-source order contrast here trades breadth for experimental control. It compares two explicitly defined pipelines at fixed settings and records exactly which order makes a particular decision fail. The association can motivate a follow-up test with actual transmission and varied gain controls. It cannot stand in for that test because the present experiment omits network loss and device processing.

Where a review decision is threshold based, a deployment report should specify the order and settings of the signal-processing chain and give error counts by class. A rank statistic can summarize detector ordering, and an error transition table can show whether a system has shifted errors to a different group. The latter is valuable even if total counts tie. For example, two methods can flag the same number of bona fide sources but require review of different recordings. No application-specific cost or calibration was measured, so a cost-optimal operating rule remains a separate problem.

The present controlled intervention uses ordinary BCE and one view per source. Its scientific contribution is the error accounting and the evidence that the observed rank benefit coexists with subgroup deterioration. Extending the claim would require a locked independent partition, then genuinely new corpora or models to measure broader transport. Confirmation of a forensic mechanism would additionally require a cue definition and a measurement that can vary independently of the detector output.

\subsection{From a validation metric to a validation workflow}
For an engineer deciding whether to adopt augmentation, the condition table distinguishes several cases. An AUC increase with reductions in both error types shows improvement at both levels measured. An AUC increase with fewer false alarms and more misses indicates that a rank gain and an operating-point trade-off coexist. A near-tie in counts can still conceal a change in the individual sources that require attention. The present results exhibit these behaviors under different conditions. A validation workflow benefits from retaining source identity through the paired analysis.

Nominal condition names can also obscure the processing intervention. A label that names noise and Opus omits whether the noise is present during coding. The operator descriptions here fix that detail and define how the noise realization is shared. This isolates sensitivity to the specified order without assuming either chain represents all communication paths. It also motivates reporting codec rates alongside order. Opus64 has different settings and support from the 24~kb/s compositions, and their contrast cannot isolate bitrate. The present evidence pertains to a classifier after specified transformations. A claim about forensic cue retention would require a separate measurement that varies independently of classifier output.

\subsection{Evidence needed to extend the claim}
The conference reviewers raised related questions about broader robustness and repeatability. The journal makes progress by separating training initialization from test-source resampling, documenting how the operating rule interacts with class-specific effects, and retaining source links through the analysis. These changes address inference and diagnosis within the registered scope. They leave broader shifts unresolved. The appropriate next empirical test is whether the same fixed models preserve the observed directional trade-offs on the locked within-release cohort. A benchmark result there would still leave the external-corpus, physical channel and multi-architecture questions open.

If the held-out ranking gain weakens or the error trade-off changes, both outcomes belong in the main results. In particular, the observed A05 and A06 development errors cannot be projected directly onto A07 through A19. The official corpus documentation relates some attack algorithms across partitions, and the five-column metadata does not identify full generator lineage. The scientific conclusion should therefore be phrased in terms of the designated partitions and conditions. Claims of genuinely unseen generator robustness require an additional design with independently verified generator provenance.

\subsection{Reporting a locked follow-up}
The interpretation of a subsequent held-out comparison should distinguish transport of the aggregate direction from transport of each conditional error pattern. If both stress ranking differences remain positive, the result would extend the observed direction to the registered evaluation partition. The fixed-threshold errors would still require separate reporting by class and condition. A change in the FAKE miss direction would be informative even in the presence of a rank benefit. The development A05 and A06 observations should remain in the paper regardless of the held-out outcome because the attack sets and source populations differ.

The primary held-out set deliberately allocates equal quotas to the registered FAKE attacks. This helps make support transparent and reduces the influence of an attack merely having more available archive records. It does not reproduce real-world spoof prevalence. Attack-stratified misses and class-conditional errors therefore remain more interpretable than an isolated accuracy figure. The same selection constraints and scores must apply to every model. Reporting the number of sources and views together will also prevent a complete set of predictions from being mistaken for an equal number of independent utterances.

The four frozen detectors provide two paired comparisons. A common pattern across the held-out models would add evidence about these fit blocks and this partition. It would leave random-initialization variability beyond the two observed seeds, language transport, and architecture transfer unresolved. Conversely, an adverse held-out contrast should be documented without changing the threshold or subset in response. Any future recalibration or modification of the augmentation bank would define a new intervention and require a new evaluation contract. This separation protects the difference between a planned test and an exploratory diagnosis after observing its results.
